% CSC302 final project report template.
% Compile with pdflatex, tectonic, or Overleaf.
\documentclass[11pt]{article}
\usepackage[margin=1in]{geometry}
\usepackage{amsmath,amssymb,amsthm}
\usepackage{booktabs,tabularx}
\usepackage{hyperref}
\usepackage{listings}
\usepackage{xcolor}

% Unicode-friendly Lean listings (pdflatex): map the symbols you use.
\lstdefinelanguage{Lean}{
  keywords={theorem,lemma,def,example,by,fun,let,have,show,from,match,with,if,then,else,
    import,namespace,end,open,variable,structure,inductive,instance,termination_by,decreasing_by},
  sensitive=true, comment=[l]{--}, morecomment=[s]{/-}{-/}, string=[b]",
}
\lstset{
  language=Lean, basicstyle=\ttfamily\small, keywordstyle=\color{blue!60!black}\bfseries,
  commentstyle=\color{gray}, columns=fullflexible, keepspaces=true, frame=single,
  breaklines=true, extendedchars=true,
  literate={ℕ}{{$\mathbb{N}$}}1 {∣}{{$\mid$}}1 {≡}{{$\equiv$}}1 {≤}{{$\le$}}1 {≥}{{$\ge$}}1
    {≠}{{$\ne$}}1 {∀}{{$\forall$}}1 {∃}{{$\exists$}}1 {→}{{$\to$}}1 {↔}{{$\leftrightarrow$}}1
    {∧}{{$\wedge$}}1 {∨}{{$\vee$}}1 {¬}{{$\neg$}}1 {⟨}{{$\langle$}}1 {⟩}{{$\rangle$}}1
    {←}{{$\leftarrow$}}1 {·}{{$\cdot$}}1
}

\newtheorem{theorem}{Theorem}
\newtheorem{lemma}[theorem]{Lemma}
\theoremstyle{definition}
\newtheorem{definition}[theorem]{Definition}

\title{Your Project Title}
\author{Your Name(s)}
\date{CSC302 Final Project, \today}

\begin{document}
\maketitle

\begin{abstract}
One paragraph: which theorem you formalized, from which starting library, and the size of
your formalization (number of lines, definitions and lemmas).
\end{abstract}

\section{Introduction}
State the theorem informally. Why is it interesting? Give a small numerical example.

\section{Paper proof}
\begin{theorem}
The precise statement.
\end{theorem}
State every lemma you use, with its proof.

\section{Lean formalization}
\subsection{Definitions}
Show each definition you introduced and explain why you chose it.

\subsection{The main theorem in Lean}
\begin{lstlisting}
theorem my_main_theorem (n : ℕ) (h : 2 ≤ n) : ∃ p, Prime p ∧ p ∣ n := by
  ...
\end{lstlisting}
Explain the statement symbol by symbol. Does it say exactly what the theorem in Section 2
says? If not, what is the difference and why?

\subsection{Key lemmas}
Walk through one or two proofs line by line.

\section{Correspondence between the paper proof and the code}
\begin{tabularx}{\textwidth}{@{}lXl@{}}
\toprule
Step & Paper proof & Lean lemma \\
\midrule
1 & Every $n \ge 2$ has a prime divisor. & \texttt{exists\_prime\_dvd} \\
2 & \dots & \dots \\
\bottomrule
\end{tabularx}

\section{What Lean made us do}
Steps the paper proof treats as obvious that needed real work in Lean; edge cases; how
statements changed because the library has no subtraction.

\section{Use of AI tools}
Tools used and for what. At least one example where the AI output was wrong or unhelpful,
and how you noticed.

\begin{thebibliography}{9}
\bibitem{nng} K.~Buzzard and J.~Eugster, \emph{The Natural Number Game}.
\bibitem{mop} H.~Macbeth, \emph{The Mechanics of Proof}.
\end{thebibliography}

\appendix
\section{Output of \texttt{\#print axioms}}
Paste the output for your main theorem.

\end{document}
